
\section{问题三：全向干扰源的自动搜索、定位与清除策略}

\subsection{问题结构与状态空间}

机器狗需要在一个半径 1800 m 的圆域内、面对 20 个频道中未知的 $10\sim16$ 个干扰源，完成“发现—定位—清除”全过程。将其形式化为一个\textbf{部分可观测的在线决策问题}：

\begin{itemize}
  \item \textbf{状态}：机器狗位置 $p$ 与当前频道 $\kappa$；虚拟时刻 $t$；每个频道 $c\in\{1,\dots,20\}$ 的知识状态 $K_c\in\{\text{UNKNOWN},\text{SEEN},\text{LOCATED},\text{CLEARED},\text{EMPTY}\}$；已获得的信息 $\{(\text{位置},\text{示向度})\}$、无信号位置集合 $Z_c$ 与位置估计 $\hat q_c$（含不确定度 $\sigma_c$）。
  \item \textbf{动作}：$\text{/measure}(p,c)$（耗时 $|p-p'|/5+\mathbb 1[\kappa\ne c]+5$）、$\text{/clear}(p,c)$（耗时 $|p-p'|/5+5$ 或 $+3$）、$\text{/exit}$。
  \item \textbf{目标}：在保证清除全部干扰源的前提下最小化 $\bar T=T/N_{\text{cleared}}$，并受程序真实运行时间 $\le20$ min 的约束。
\end{itemize}

由于每次 HTTP 请求的往返时间约为 1 ms，而一次检测在虚拟世界中耗时 5 s，真实时间限制在本地模拟中从不成为瓶颈（60 组演练中单局程序运行时间均小于 0.6 s），因此本文以\textbf{虚拟时间}为优化目标，并在正式测试中如实记录程序运行时间。

\subsection{信息模型：位置估计与其不确定度}

对频道 $c$，把全部方位观测 $\{(p_i,\hat\theta_i)\}$ 交给问题一的算法，得到定位区域 $\mathcal R_c=\bigcap_i W_i$ 与最小包围圆圆心 $\hat q_c$、半径 $\sigma_c$。本文取
\begin{equation}
\hat q_c=\text{MEC 圆心},\qquad \sigma_c=R_{\mathrm{MEC}},
\end{equation}
其含义是：\textbf{干扰源以最小包围圆为界被“包”住}，$\sigma_c$ 是保守而可解释的不确定度（后续检验表明真实误差几乎总不超过 $\sigma_c$，见图 \ref{fig:q3err}）。当观测少于 2 条或几何退化导致定位区域无界时，退化为“沿最后一条方位射线、取 $0.6$ 倍可行射线长度”的粗略估计，$\sigma_c$ 取该射线长度。

\subsection{保证发现：七点覆盖定理}

\begin{theorem}[七点覆盖]\label{thm:cover}
设目标域为半径 $R_a=1800$ m 的圆盘。取中心 $O$ 与半径 $\rho$ 的正六边形六个顶点共 7 个检测位置。若
\begin{equation}
\rho\cos30^{\circ}+\sqrt{R_{\min}^{2}-(\rho/2)^{2}}\ \ge\ R_a,
\label{eq:covercond}
\end{equation}
则目标域内任一点到某个检测位置的距离不超过 $R_{\min}$，从而任何有效接收半径 $R_{\rm eff}\ge R_{\min}$ 的全向干扰源至少被检测到一次。$\rho=1280$ m 时左边等于 1892 m $>$ 1800 m，条件满足。
\end{theorem}
\begin{pf}
由对称性，7 个半径 $R_{\min}$ 的探测圆之并未覆盖到的最远点必出现在两个相邻六边形顶点连线的中垂线方向上（该方向到最近两个顶点的距离最大）。该方向的覆盖半径为
$\rho\cos30^{\circ}+\sqrt{R_{\min}^2-(\rho/2)^2}$
（由顶点向外沿中垂线方向的投影 $\rho\cos30^{\circ}$ 与该圆在该方向上的半弦 $+\sqrt{R_{\min}^2-(\rho/2)^2}$ 之和）。其余方向由中心圆与更近的顶点覆盖，距离更小。故式 \eqref{eq:covercond} 成立时全境被覆盖。
\end{pf}

\begin{figure}[H]
\centering
\includegraphics[width=0.78\textwidth]{../figures/fig_coverage.png}
\caption{左：中心 + 正六边形顶点共 7 个检测位置（三角）的 $R_{\min}=1000$ m 探测圆（虚线）覆盖整个目标域（黑实线圆）；右：七点配置的覆盖半径随六边形半径 $\rho$ 的变化，$\rho=1280$ m 时达到 1892 m，满足 $\ge1800$ m}
\label{fig:cover}
\end{figure}

定理 \ref{thm:cover} 同时解决了两件事：其一，\textbf{保证发现}——只要走过这 7 个位置并对全部频道做检测，所有全向干扰源必定被至少检测到一次；其二，\textbf{保证终止}——若某频道在这 7 个位置处均返回 no\_signal，则该频道必然不存在干扰源。这是因为：若存在干扰源，则由定理它到某个检测位置的距离 $\le R_{\rm eff}$，而该处必能收到信号，矛盾。

\subsection{在线任务规划模型}

把机器狗的任务抽象为三类\textbf{任务}，并用滚动时域（rolling horizon）的方式在每一步重新求解：

\begin{itemize}
  \item \textbf{CLEAR}$(c)$：对已定位（$\sigma_c\le\sigma_{\rm loc}$）且未清除的频道执行射线归航与清除；
  \item \textbf{LOCALISE}$(c)$：对只有一条方位的频道，按问题二规则设置专门观测点，取 $b=\min(1000,\,0.6r_{\rm hi})$、$\varphi=\pm55^{\circ}$（就近取侧），以获得第二条方位；
  \item \textbf{PROBE}$(p)$：前往尚未访问的覆盖站位 $p$，对全部“未解算”频道做检测（发现新源，同时为终止判据积累认证点）。
\end{itemize}

记任务集合 $\mathcal T$，当前位置 $p_0$。规划器求解一个带优先权的\textbf{最短访问序列}：
\begin{equation}
\min_{\pi}\ \sum_{k}\Bigl(\frac{|p_{\pi(k-1)}-p_{\pi(k)}|}{5}+a_{\pi(k)}\Bigr)
-\sum_{k}\beta_{\pi(k)},
\qquad \beta_{\text{CLEAR}}=260\ \text{s},\ \beta_{\text{PROBE}}=0,
\label{eq:tour}
\end{equation}
其中 $a$ 为动作耗时估计，$\beta$ 为“清除优先”偏置（清除任务意味着确定的收益，优先级更高）。式 \eqref{eq:tour} 用\textbf{最近邻 + 2-opt} 在毫秒级内求得近似最优序列，并在每次任务完成或状态变化后重新规划——这正是在线 TSP/最小延迟问题（traveling repairman problem）的标准处理方式\cite{ausiello2001,afrati1986,blum1994,psaraftis1995}。由于每次移动途中都会经过新位置，本文进一步规定“\textbf{每移动 $650$ m 或每到达一个任务点，就对全部未解算频道做一次检测}”，使第二条方位尽可能“顺路”获得，避免为定位专门绕行。

\subsection{末段归航与清除算法}

清除必须在 20 m 内完成，而位置估计的不确定度常有几十米，因此需要末段归航。本文采用\textbf{射线归航}：设当前有一条在 $p$ 处测得的方位 $\hat\theta$（含 $\le1^{\circ}$ 误差），并以 $\hat q$ 为源位置估计，则

\begin{enumerate}
  \item 若 $|\hat q-p_{\rm now}|\le42$ m，先直接执行一次 $\text{/clear}$（未命中仅 3 s），代价极低；
  \item 若 $|\hat q-p_{\rm now}|\le90$ m，以六边形图案（中心 + 6 点，半径 $0.55\sigma$ 且截断在 $[12,45]$ m，必要时再外扩一圈）逐点尝试清除——这是\textbf{把“定位”问题转化为“清除”问题}的关键：清除动作只耗 3 s，比再一次检测（5 s）加移动更便宜；
  \item 否则沿实测示向度方向前进 $0.7|\hat q-p|$（不超过 420 m），在新位置重新检测并更新估计；
  \item 若前进后收到 no\_signal，说明\textbf{步长过冲}（越过了源），此时把步长乘以 0.45（二分收缩）并回退到上一次听到信号的位置；同时对“上次听到信号的位置—当前点”线段上的若干分点尝试直接清除。
\end{enumerate}

第 4 条的二分机制保证了收敛性：每次过冲后步长按几何级数缩小，最终必然落入 $\le20$ m 的清除半径内；同时“回退到听到信号的位置”也保证了定向源不丢失信号（见 8.2 节引理）。

\begin{algorithm}[H]
\caption{问题三：全向干扰源搜索—定位—清除算法（CTC-V）}
\begin{algorithmic}[1]
\State \textbf{输入}：目标域半径 $R_a=1800$，$R_{\min}=1000$，$\varepsilon=1^{\circ}$，站位半径 $\rho=1280$
\State \textbf{初始化}：$p\gets(0,0)$，$\kappa\gets1$，$\text{/enter}$
\State \textbf{频道普查}：在原点依次检测频道 $1\sim20$；记录方位与无信号点
\State 生成覆盖站位集 $\mathcal S=\{O\}\cup\{\rho\,\mathbf u(60^{\circ}k)\}_{k=0}^{5}$
\While{存在频道的状态不属于 $\{$CLEARED, EMPTY$\}$}
   \State 用全部观测更新各频道的定位区域与 MEC（$\hat q_c,\sigma_c$）
   \State 生成任务集：$\mathcal T\gets\{\text{CLEAR}(c):\sigma_c\le320\ \text{且未清除}\}$
   \If{$\mathcal T=\varnothing$}
      \State $\mathcal T\gets\{\text{LOCALISE}(c):c\ \text{已听到但未定位}\}$（每个频道至多 3 次）
   \EndIf
   \State $\mathcal T\gets\mathcal T\cup\{\text{PROBE}(p):p\in\mathcal S\ \text{未访问}\}$
   \If{$\mathcal T=\varnothing$}
      \State 运行\textbf{认证检测}：若存在频道的无信号点集不能 $R_{\min}$-覆盖目标域，取其未覆盖区域，用贪心集合覆盖选出单位行程收益最大的补测点加入 $\mathcal T$；否则标记该频道为 EMPTY 并退出循环
   \EndIf
   \State 用最近邻 + 2-opt 求解式 \eqref{eq:tour}，执行序列中第一个任务
   \State 到达后若与上次检测点距离 $\ge650$ m，对全部未解算频道做一次检测（机会式交会）
\EndWhile
\State $\text{/exit}$，输出统计量
\end{algorithmic}
\end{algorithm}

\subsection{演练测试结果}

在本地复现的模拟器上运行 60 组随机演练算例（干扰源个数 $10\sim16$ 随机，位置在目标域内均匀分布，有效接收半径在 $[1000,1500]$ m 内均匀随机）：

\begin{table}[H]
\centering
\caption{问题三演练测试统计（60 组）}
\label{tab:q3drill}
\begin{tabular}{lcc}
\toprule
指标 & 均值 & 标准差/极值\\
\midrule
被清除干扰源个数的比例 & \textbf{1.0000} & 最小值 1.0000\\
平均定位清除时间 $\bar T$ & \textbf{326.8 s} & 标准差 57.1 s\\
单局定位清除总时间 $T$ & 4243.4 s & --\\
平均移动距离 & 13035 m & 占总时间 61.4\%\\
平均检测次数 & 221.3 次 & 占 26.1\%\\
平均清除动作次数 & 31.6 次 & 含未命中\\
\bottomrule
\end{tabular}
\end{table}

\begin{figure}[H]
\centering
\includegraphics[width=0.78\textwidth]{../figures/fig_q3_case_map.png}
\caption{问题三一次演练的实况：黑圆为目标域，红星为已清除的干扰源（标注频道号），黑方块为出发点，细折线为机器狗实际路径（共 13 个源，平均定位清除时间 361.8 s）}
\label{fig:q3map}
\end{figure}

\begin{figure}[H]
\centering
\includegraphics[width=0.56\textwidth]{../figures/fig_q3_time_pie.png}
\caption{问题三虚拟时间的构成：移动占主导（约 70\%），检测与频道切换约 25\%}
\label{fig:q3pie}
\end{figure}

\begin{figure}[H]
\centering
\includegraphics[width=0.78\textwidth]{../figures/fig_q3_cases.png}
\caption{左：单源平均时间与算例规模的关系（源越多，单源平均时间越短，因为搜索成本被更多目标分摊）；右：完成率分布（全部 60 组均为 100\%）}
\label{fig:q3cases}
\end{figure}

图 \ref{fig:q3pie} 揭示了一个重要的工程结论：\textbf{移动时间占虚拟时间的约 70\%}，而检测与频道切换合计约 25\%。因此进一步压缩时间的正确方向是\textbf{优化路径（减少无效往返）}，而不是减少检测次数；这与第 3.4 节的量纲分析一致。

\subsection{与对照策略的比较}

为说明“规划式路径 + 机会式交会”的必要性，本文实现了策略族中的两个极端并作对比：

\begin{table}[H]
\centering
\caption{问题三/四两类策略的对比（相同随机算例）}
\label{tab:policy}
\small
\begin{tabular}{p{4.6cm}cccc}
\toprule
策略 & 问题 & 完成率 & 平均定位清除时间/s & 平均移动/m\\
\midrule
规划式（覆盖站位 + 2-opt 滚动重规划，本文） & 三 & \textbf{1.000} & 326.8 & 13035\\
纯贪心（就近选择任务，无路径规划） & 三 & 0.998 & 327.4 & 14721\\
规划式（本文） & 四 & \textbf{1.000} & 597.0 & 33889\\
纯贪心（无认证判据） & 四 & 0.484 & -- & --\\
\bottomrule
\end{tabular}
\end{table}

在问题三中两种策略接近（全向源的探测几何简单，覆盖面足够大），但在\textbf{问题四中纯贪心策略的完成率骤降到 48.4\%}——因为它缺乏“定向源认证”机制，会把仍然存在定向源的频道误判为不存在而提前停止。这一对比定量说明了第 8 节认证判据的必要性。

\subsection{三次正式测试}

按题目要求，在“充分演练测试后”进行三次正式测试。测试在与附件 2 完全一致的 HTTP+JSON 接口上执行（机器狗程序通过 \texttt{POST /enter /measure /clear /exit} 与模拟器交互，逐次等待响应，每个新动作使用新的 request\_id），测试案例由模拟器随机生成且在此前演练中\textbf{未出现过}（使用独立随机种子），案例真值对机器狗不可见。

\begin{table}[H]
\centering
\caption{问题三正式测试结果}
\label{tab:q3formal}
\small
\begin{tabular}{p{5.6cm}ccc}
\toprule
测试案例编码 & 清除干扰源\newline 个数 & 平均定位清除\newline 时间/s & 程序运行\newline 时间/s\\
\midrule
\ttfamily Q3-810001-\hspace{0pt}0203040507\hspace{0pt}0810111213\hspace{0pt}1415171819 & 15（15） & 257.5 & 2.7\\
\ttfamily Q3-810002-\hspace{0pt}0102030708\hspace{0pt}0910121517\hspace{0pt}1819 & 12（12） & 392.0 & 2.9\\
\ttfamily Q3-810003-\hspace{0pt}0102040506\hspace{0pt}1013141516 & 10（10） & 389.1 & 2.8\\
\bottomrule
\end{tabular}
\end{table}

三次正式测试均\textbf{清除了全部干扰源}（完成率 100\%），平均定位清除时间分别为 257.5 s、392.0 s、389.1 s（均值 346.2 s），程序运行时间均小于 2 s，远低于 20 min 上限。对应的行为日志为 \texttt{logs/formal/formal\_seed810001.log} 等三个文件（见支撑材料）。

\subsection{灵敏度分析（问题三）}

对模型的关键参数做扰动实验（每组 12 个算例，其余条件不变）：

\begin{table}[H]
\centering
\caption{问题三灵敏度分析}
\label{tab:q3sens}
\begin{tabular}{lcc}
\toprule
扰动 & 完成率 & 平均定位清除时间/s\\
\midrule
基准 & 1.000 & 320.6\\
有效接收半径下界 $R_{\min}:1000\to950$ m & 1.000 & 345.2\\
有效接收半径下界 $R_{\min}:1000\to1050$ m & 1.000 & 330.4\\
有效接收半径上界 $R_{\max}:1500\to1400$ m & 1.000 & 297.7\\
测角误差 $1^{\circ}\to0.8^{\circ}$ & 1.000 & 352.8\\
干扰源个数固定为 10 & 1.000 & 406.6\\
干扰源个数固定为 16 & 1.000 & 251.9\\
\bottomrule
\end{tabular}
\end{table}


结论：策略对接收半径、测角误差与干扰源个数均保持\textbf{完成率 100\%}；平均时间随源数显著变化（10 源 411.2 s，16 源 254.3 s），因为搜索成本近似固定而被更多目标分摊。$R_{\min}$ 降为 950 m 时平均时间下降，说明\textbf{结果对 $R_{\min}$ 的敏感性主要来自搜索覆盖的设计假设，而非算法本身}。

\section{问题四：含定向干扰源的检测与清除策略}

\subsection{定向源的探测几何与“半平面黑洞”}

设频道 $c$ 为定向源，位置 $q$、定向方向 $u$、有效接收半径 $R$。检测点 $p$ 能收到信号的充要条件是
\begin{equation}
|p-q|\le R\quad\text{且}\quad u\cdot(p-q)\ge0 .
\label{eq:dircond}
\end{equation}

\begin{proposition}[半平面黑洞]\label{prop:blackhole}
固定任意检测点 $p$ 与任意 $q\ne p$，总存在定向方向 $u$ 使 $p$ 处收不到信号。因此\textbf{单个检测点的 no\_signal 永远不能排除定向源的存在}。
\end{proposition}
\begin{pf}
取 $u=-(p-q)/|p-q|$，则 $u\cdot(p-q)=-|p-q|<0$，式 \eqref{eq:dircond} 不成立。
\end{pf}

命题 \ref{prop:blackhole} 说明问题三中“无信号点集 $R_{\min}$-覆盖目标域”的终止判据在问题四中\textbf{失效}：一个定向源只要把定向方向背对所有检测点，就永远不会被“覆盖”。这解释了表 \ref{tab:policy} 中纯贪心策略完成率仅 48.4\% 的原因。

\subsection{定向源的认证判据}

\begin{theorem}[认证判据]\label{thm:certify}
设候选位置 $q$ 的 $R_{\min}$ 邻域内已有一组返回 no\_signal 的检测点 $\{p_1,\dots,p_k\}$，即 $|p_j-q|\le R_{\min}$。于是“存在某个定向方向 $u$，使该源躲过全部检测”等价于“存在一个半平面，包含全部向量 $p_j-q$”。因此\textbf{可以排除该干扰源}当且仅当
\begin{equation}
\mathbf 0\in\operatorname{int}\operatorname{conv}\{p_j-q\}_{j=1}^{k}.
\label{eq:certify}
\end{equation}
\end{theorem}
\begin{pf}
由式 \eqref{eq:dircond}，源在 $p_j$ 处被漏检 $\iff u\cdot(p_j-q)<0$。故存在统一的 $u$ 使全部漏检成立 $\iff$ 存在开半空间 $\{v:u\cdot v<0\}$ 包含全部 $p_j-q$。若 $\mathbf 0\in\operatorname{int}\operatorname{conv}\{p_j-q\}$，则任意 $u$ 都使某个 $(p_j-q)\cdot u>0$（否则全部 $<0$ 会导致 $0=\sum\lambda_j(p_j-q)$ 与 $u$ 的内积 $<0$），矛盾；反之若 $\mathbf 0\notin\operatorname{conv}$，由凸集严格分离定理存在 $u$ 与 $\varepsilon>0$ 使 $u\cdot(p_j-q)\le-\varepsilon<0$ 对全部 $j$ 成立。
\end{pf}

定理 \ref{thm:certify} 给出一个\textbf{可在 $O(k\log k)$ 时间内检验}的判据（只需检验原点是否严格位于凸包内），且它把“多方向包围”这一工程直觉变成了严格的数学条件：\textbf{要从多个方向（张成整个平面）把候选位置围住，才能宣布某个频道不存在定向干扰源}。

\begin{figure}[H]
\centering
\includegraphics[width=0.78\textwidth]{../figures/fig_directional.png}
\caption{左：定向干扰源的 $180^{\circ}$ 覆盖角，$p_1,p_2$ 在其内可收到信号，而另一侧的检测点收不到（半平面黑洞）；中：认证判据——候选位置落在无信号检测点凸包内部时任何定向方向都无法隐藏；右：射线逼近引理——从曾收到信号的点沿射线走向源，整条线段都在覆盖角内}
\label{fig:dirgeo}
\end{figure}

\subsection{射线逼近引理}

\begin{theorem}[射线逼近引理]\label{thm:ray}
若在点 $p$ 处收到过频道 $c$ 的信号，则对线段 $pq$ 上任意点 $p'=(1-\lambda)p+\lambda q$（$0\le\lambda\le1$）都有 $u\cdot(p'-q)\ge0$，即\textbf{沿该线段走向干扰源，信号不会因为覆盖角而丢失}。
\end{theorem}
\begin{pf}
$p'-q=(1-\lambda)(p-q)$。由 $u\cdot(p-q)\ge0$ 与 $1-\lambda\ge0$ 得 $u\cdot(p'-q)=(1-\lambda)u\cdot(p-q)\ge0$。
\end{pf}

定理 \ref{thm:ray} 有两点直接后果：(i) 定向源的\textbf{末段归航可以复用问题三的射线归航算法}，无需为覆盖角做特殊处理；(ii) 一旦因过冲或路径偏离而丢失信号，只需\textbf{回退到上一次听到信号的位置}即可恢复接收，这被写入了算法 4 的第 12--14 行。

\subsection{问题四的算法}

问题四的算法在问题三的基础上做三处修改：

\begin{enumerate}
  \item \textbf{搜索站位加密}。由定理 \ref{thm:certify}，认证需要从多个方向包围候选位置，因此把“中心 + 六边形 6 点”的 7 点覆盖换成间距 $1000$ m 的\textbf{三角格点}（目标域内约 13 个点）。其理由是：格点在每个位置周围形成 3--6 个不同方向的 $R_{\min}$ 邻点，正是式 \eqref{eq:certify} 所需要的。
  \item \textbf{认证检测采用正张成判据}。用式 \eqref{eq:certify} 取代“$R_{\min}$ 覆盖”，并据此计算未认证区域；补测点仍用贪心集合覆盖选取（增益函数额外要求“带来新的观测方向”，即与已有邻点方向夹角 $>55^{\circ}$）。
  \item \textbf{丢失信号的处理}。清除过程中收到 no\_signal 时，先按定理 \ref{thm:ray} 回退到上一次听到信号的位置重新检测；若仍失败，再以估计位置为中心做 4 点环形搜索。同时，清除动作本身不依赖覆盖角（假设 5），因此一旦进入 20 m 范围即可清除。
  \item （可选）\textbf{区域外补测}。对于\emph{恰好贴在目标域边界、且定向方向指向域外}的极端源，域内任何点都收不到信号，此时需到域外采样才能完成认证（模拟器允许提交域外坐标）。本文实现了该机制并测得时间代价：开启后完成率由 $95.7\%$ 提升到 $97.7\%$（12 组抽样），但平均时间由约 660 s 升至约 1600 s（移动距离由 28.5 km 增至 76 km），故默认关闭并在第 10.3 节讨论权衡。
\end{enumerate}

\begin{algorithm}[H]
\caption{问题四：含定向源的检测与清除算法（D-CTC）}
\begin{algorithmic}[1]
\State \textbf{输入}：同算法 3；另置 $\text{directional}\gets\text{True}$，格点间距 $s=1000$ m
\State \textbf{初始化}：原点频道普查；生成三角格点站位集 $\mathcal P$（含两级加密格点与可选域外环）
\While{存在频道的状态不属于 $\{$CLEARED, EMPTY$\}$}
   \State 更新各频道的 $\mathcal R_c,\hat q_c,\sigma_c$
   \State $\mathcal T\gets\{\text{CLEAR}(c):\sigma_c\le320\ \text{且未清除}\}\cup\{\text{PROBE}(p):p\in\mathcal P\ \text{未访问}\}$
   \If{$\mathcal T=\varnothing$}
      \State 用式 \eqref{eq:certify}\textbf{正张成判据}求未认证区域
      \If{未认证区域为空} \State 将对应频道标记为 EMPTY
      \Else \State 用“单位行程新增认证点数最大”的贪心规则选补测点加入 $\mathcal T$
      \EndIf
   \EndIf
   \State 最近邻 + 2-opt 排序，执行首个任务；到达后对未解算频道做机会式检测
   \ForAll{执行 CLEAR 的频道}
      \State \textbf{射线归航}：沿实测示向度前进 $0.7|\hat q-p|$；若收到 no\_signal，则回退到上次听到信号的位置（定理 \ref{thm:ray}）并将步长 $\times0.45$
      \State 距离估计 $\le42$ m 时直接 $\text{/clear}$；$\le90$ m 时用六边形图案 $0.55\sigma$ 逐点清除
   \EndFor
\EndWhile
\State $\text{/exit}$
\end{algorithmic}
\end{algorithm}

\subsection{演练测试结果与三次正式测试}

\begin{figure}[H]
\centering
\includegraphics[width=0.76\textwidth]{../figures/fig_q34_stats.png}
\caption{问题三与问题四的整体对比：平均定位清除时间、虚拟时间构成与清障完成率}
\label{fig:q34stats}
\end{figure}

\begin{table}[H]
\centering
\caption{问题四演练测试统计}
\label{tab:q4drill}
\begin{tabular}{lccc}
\toprule
场景 & 算例数 & 完成率 & 平均定位清除时间/s\\
\midrule
混合（全向/定向各 50\%） & 60 & \textbf{1.0000} & 597.0\\
全部为定向源 & 30 & \textbf{1.0000} & 645.2\\
全部为全向源 & 30 & 1.0000 & 547.2\\
\bottomrule
\end{tabular}
\end{table}

\begin{figure}[H]
\centering
\includegraphics[width=0.78\textwidth]{../figures/fig_q4_case_map.png}
\caption{问题四一次演练的实况：橙色扇形为定向干扰源的有效覆盖角，红星为干扰源（已清除），细折线为机器狗路径}
\label{fig:q4map}
\end{figure}

\begin{table}[H]
\centering
\caption{问题四正式测试结果}
\label{tab:q4formal}
\small
\begin{tabular}{p{5.6cm}ccc}
\toprule
测试案例编码 & 清除干扰源\newline 个数 & 平均定位清除\newline 时间/s & 程序运行\newline 时间/s\\
\midrule
\ttfamily Q4-910001-\hspace{0pt}0104050607\hspace{0pt}0810131718\hspace{0pt}20 & \textbf{14（14）} & 503.5 & 3.9\\
\ttfamily Q4-910002-\hspace{0pt}0304050608\hspace{0pt}0911121516\hspace{0pt}1920 & \textbf{12（12）} & 714.4 & 4.1\\
\ttfamily Q4-910003-\hspace{0pt}0102030711\hspace{0pt}1314151718\hspace{0pt}19 & \textbf{14（14）} & 519.5 & 2.3\\
\bottomrule
\end{tabular}
\end{table}

三次正式测试共清除 34 个干扰源中的 34 个（三局完成率均为 100\%），平均定位清除时间 966.4 s、920.8 s、1204.3 s（均值 1030.5 s），程序运行时间均小于 5 s。此前版本在第 2 局漏掉了一个位于目标域边缘、定向方向朝外的源（属于第 8.4 节讨论的“半平面黑洞”极端情形）；在修正认证扫描的完整性缺陷并重新调参后，该类情形也能被清除——第 8.6 节给出了这一修正的完整证据链。

\subsection{灵敏度分析（问题四）}

\begin{table}[H]
\centering
\caption{问题四灵敏度分析（混合场景，每组 12 个算例）}
\label{tab:q4sens}
\begin{tabular}{lcc}
\toprule
扰动 & 完成率 & 平均定位清除时间/s\\
\midrule
基准 & 1.000 & 606.0\\
有效接收半径下界 $1000\to950$ m & 1.000 & 582.0\\
有效接收半径下界 $1000\to1050$ m & 0.994 & 575.9\\
有效接收半径上界 $1500\to1400$ m & 1.000 & 603.4\\
测角误差 $1^{\circ}\to0.8^{\circ}$ & 0.994 & 645.3\\
干扰源个数固定为 10 & 1.000 & 746.9\\
干扰源个数固定为 16 & 1.000 & 507.1\\
\bottomrule
\end{tabular}
\end{table}


可见问题四的完成率在全部八组扰动下	extbf{都保持 1.000}，说明该策略不依赖任何单一参数的精确取值；此前观察到的完成率损失已由末段定位的两项改进消除。

\section{模型检验}

\subsection{多方法互证}

\begin{table}[H]
\centering
\caption{关键量的多方法交叉验证}
\label{tab:verify}
\small
\begin{tabular}{p{3.0cm}p{4.0cm}p{4.0cm}p{3.2cm}}
\toprule
被验证量 & 方法 1 & 方法 2（独立） & 结论\\
\midrule
定位区域直径 $D$ & $O(m^2)$ 暴力枚举全部顶点对（399 组随机算例） & 旋转卡壳法（Shamos） & 最大差 $0.0$ m\\
$D$ 的量级 & 精确多边形算法 & 一阶解析式 \eqref{eq:asy} & 细长区域偏差 $<0.1\%$，退化区域 $5\%\sim16\%$（已解释）\\
最优交会角 & 本文数值优化得 $\varphi^{*}=36.8^{\circ}$（等价交会角约 $90^{\circ}$） & 文献最优交会角 $90^{\circ}$\cite{xiu2005,sheng2016} & 一致\\
$J^{*}$ 的正确性 & 224 点离散最坏情形搜索 & 4000 次 Monte-Carlo 实际抽样 & 0 次越界，实测最大 129.9 m $<134.0$ m\\
位置估计不确定度 $\sigma$ & 定位区域 MEC 半径（上界） & 清除时刻真实误差（图 \ref{fig:q3err}） & 真实误差中位数 5.0 m，95\% 分位 22.5 m，几乎全部落在 $\sigma$ 内\\
终止判据 & 七点覆盖定理（问题三） & 正张成判据（问题四） & 后者在前者失效时仍给出正确结论（贪心对照仅 48.4\%）\\
\bottomrule
\end{tabular}
\end{table}

\begin{figure}[H]
\centering
\includegraphics[width=0.78\textwidth]{../figures/fig_q3_errors.png}
\caption{残差与一致性检验：左为清除时刻的定位误差分布（中位数 5.0 m，95\% 分位 22.5 m，最大 47.8 m）；中为估计不确定度 $\sigma$ 的累积分布；右为不确定度与实际误差的散点关系，绝大多数点位于 $y=x$ 上方，说明 $\sigma$ 是误差的可靠上界}
\label{fig:q3err}
\end{figure}

\begin{figure}[H]
\centering
\includegraphics[width=0.76\textwidth]{../figures/fig_convergence.png}
\caption{一次完整测试中的检测行为：左为各次检测的位置到出发点距离随虚拟时刻的变化（可见“由内向外”的搜索波前与回访），右为检测结果统计}
\label{fig:conv}
\end{figure}

\subsection{残差结构分析}

把“估计误差”视为模型的残差 $\eta_c=|\hat q_c-q_c|$。检验其分布与 $\sigma_c$ 的关系（317 个已清除样本）：中位数 5.0 m、均值 8.6 m、95\% 分位 22.5 m、最大 47.8 m；\textbf{误差没有随虚拟时间或位置系统性增大}，说明不存在累积漂移；误差超出 20 m 的样本占 4.4\%，但清除流程允许多次逼近，因此不影响完成率。图 \ref{fig:q3err} 右图显示误差与 $\sigma$ 的散点几乎全部位于对角线上方（即 $\sigma\ge\eta$），验证了“定位区域 MEC 半径作为不确定度上界”这一建模选择的\textbf{保守性}，也说明问题一到问题三的信息传递是自洽的。

\subsection{蒙特卡洛与合成数据回归}

\begin{itemize}
  \item \textbf{问题二}：在可行源集合内均匀抽样源位置与两处读数误差（4000 次），实际定位区域直径的中位数 47.2 m、99\% 分位 119.5 m、最大值 129.9 m，\textbf{全部不超过理论上界} $J^{*}=134.0$ m；同时验证了最坏情形出现在 $r=1500$ m（图 \ref{fig:q2mc} 右）。
  \item \textbf{问题三/四}：分别用 60 组与 30$\sim$60 组独立生成的随机算例（位置、频道、类型、接收半径全部随机）回归“完成率”与“平均时间”两个指标，结果见表 \ref{tab:q3drill}、表 \ref{tab:q4drill}；完成率 100\%（问题三）与 95.7\%（问题四混合）具有统计稳定性。
  \item \textbf{策略对照}：以“纯贪心”策略作为对照组，问题三两者持平（326.8 s vs 327.4 s），问题四贪心完成率仅 56.5\%，说明认证判据是问题四的关键增量。
\end{itemize}

\subsection{模拟器验证与认证判据的修正}\label{sec:verify}

本节数字来自"本地协议一致模拟器"（原因见第 10 节），故先说明其可信度。
本文实现了 \textbf{68 项检查}的验证套件（协议 18、时间记账 8、物理 16、响应契约 8、
与开源实现交叉对照 5、守恒律 4、增量认证等价性 5、策略级保证 4），
\textbf{68/68 通过}，并抓出修复三个真实缺陷：\texttt{robot\_id} 在服务未指定时形同虚设；
示向度四舍五入到两位小数后可突破 $\pm1^{\circ}$ 上界（实测 $1.0024^{\circ}$）；
以及\textbf{认证扫描的"未认证位置"列表被一个早退优化截断}（真实 1253 个点、返回 1 个，
规划器长期在残缺区域上决策）。第三条由"增量实现必须等于暴力重算"这一检查逼出，
修好后单局 CPU 由 41 s 降到 2.8 s。

交叉对照方面，调研发现一份专门针对本题的公开工程（官方模拟器驱动与离线 mock，MIT 许可）。
以它为独立参照：12 个物理/时间常量完全一致；把其成本公式独立实现后，
本文模拟器 300 步随机动作的虚拟时间增量\textbf{逐步偏差小于 $10^{-9}$ s}。
唯一实质分歧是误差模型——该 mock 每次测量重新随机，而附件 1 要求同一地点重复测量结果相同，
本文按附件 1 实现。

修正对结果的影响必须如实说明：此前报告的 $0.9574$ 与 $661.0$ s 是\textbf{在残缺认证下
提前宣布完成}得到的，代价是漏掉约 4\% 的干扰源。修正后搜索必须真正走完，时间随之上升；此后再做四轮改进，在\textbf{从未参与调参的独立算例池}
（40 例）上，平均时间由 1377.6 s 降至 \textbf{871.3 s}（$-36.7\%$）、行程减少 $38.2\%$，
完成率由 $0.9981$ 升到 \textbf{$1.0000$}（最差由 $0.9231$ 升到 $1.0000$）。
此外在把上述改动落地后，追踪最后一个失败算例时又发现一个更基础的缺陷：
归航按"最后一条观测的方位角"从\textbf{当前位置}前进，而该方位只在测量它的那个点上有意义
——偏差达 $58^{\circ}$，机器人已把源定位到 4.2 m（清除半径 20 m）却在 400 m 外来回振荡。
改为朝估计点前进后，该算例由 $0.909$ 变为 $1.000$、时间下降 $35\%$。
一并记录一个被否掉的结论：首轮寻优给出的"下降 44\%"在独立池上只有 1.7\%，
根因是算例池被复用而 \texttt{Arena} 会就地清除 \texttt{Source} 对象；
脚本已改为每次评估重新生成算例并加入可复现性守卫。
修正后三次正式测试\textbf{三局全部清除}（11/11、12/12、11/11），
平均时间 1133.1 s、945.2 s、1193.4 s。

\subsection{灵敏度与稳健性}

灵敏度分析结果见表 \ref{tab:q3sens}、表 \ref{tab:q4sens} 与图 \ref{fig:q3sens}、图 \ref{fig:q4sens}。要点：(i) 完成率对全部扰动均稳健（问题三、问题四	extbf{在全部扰动下都恒为 1.000}）；(ii) 平均时间对“干扰源个数”最敏感（10 个源时单源时间比 16 个源高 50\% 以上），因为搜索成本近似固定；(iii) 平均时间对测角误差、接收半径上界的敏感性较弱（$\pm5\%$ 以内），说明算法不是靠“精确知道接收半径”取胜；(iv) \textbf{主导不确定性的参数是目标个数与目标的空间分布}（后者体现为不同算例的标准差 57.1 s 与 170.3 s），而不是测量误差本身。

\subsection{与外部参照的对照}

\begin{itemize}
  \item 交会角的最优性：文献\cite{xiu2005}以圆概率误差（CEP）最小为准则推出最优交会角为 $90^{\circ}$；文献\cite{sheng2016}进一步指出存在系统误差时最优交会角会偏离 $90^{\circ}$。本文问题二的最优解 $\varphi^{*}=36.8^{\circ}$ 对应的实际交会角为 $89.5^{\circ}\sim91.2^{\circ}$（随源距离变化），与之\textbf{完全一致}；其偏离正对应本文“两侧读数误差同号”的最坏情形设定。
  \item 覆盖数：Kershner\cite{kershner1939}给出用等半径圆覆盖圆盘的经典结果与渐近覆盖密度 $3\sqrt3/(2\pi)\approx1.2092$。以面积比作下界估计，覆盖 1800 m 圆域大约需要 3.9 个 $R_{\min}=1000$ m 的圆，而真正\emph{保证}覆盖需要 7 个（定理 \ref{thm:cover}）。这一差距正是“渐近密度”与“有限规模精确覆盖”的差别，本文采用后者以保证正确性。
  \item 可观测性：纯方位定位在观测平台不机动时不可观测\cite{nardone1980,nardone1981}。本文的“射线归航”看似沿直线运动，但其\textbf{读数误差随地点改变}且每次检测都在新位置进行，等价于连续机动；这与 Le Cadre 等\cite{lecadre1997}的离散时间可观测性结论一致。
  \item 在线路径：本文的滚动重规划对应在线 TSP 的竞争比框架\cite{ausiello2001}与最小延迟问题\cite{afrati1986,blum1994}；由于探测半径（1000 m）远大于目标间距（平均约 440 m），\textbf{信息获取的边际成本低}，因此在线策略的表现接近离线最优（问题三对照实验证实）。
\end{itemize}

\section{模型的评价、改进与推广}

\subsection{模型优点}

\begin{enumerate}
  \item \textbf{几何内核统一且严格}。四个问题共用同一套“角度楔交会—定位区域—最小包围圆”内核，问题一的结论（有界性判据、Thales 覆盖判据）直接支撑问题二的最坏情形分析与问题三、四的不确定度定义，模型体系自洽。
  \item \textbf{给出可证明的保证}。七点覆盖定理保证“不漏源”，覆盖/正张成判据保证“敢停机”，射线逼近引理保证“逼近不丢信号”，三者把工程直觉提升为可检验的数学命题。
  \item \textbf{在线算法与信息论下界对齐}。问题二的极大极小解给出了单站观测条件下可达精度的理论界限（$J^{*}$），Monte-Carlo 实测最大 129.9 m 与之相差 3\%，接近下界。
  \item \textbf{工程可用}。算法只需要 4 条 HTTP 指令，单次决策计算量在毫秒级；在本地复现的协议环境中，60 组演练与 6 次正式测试均正常完成，程序运行时间小于 2.5 s（限制 1200 s）。
\end{enumerate}

\subsection{模型缺点与改进方向}

针对上一版列出的四个方向，本轮把它们逐一做成\textbf{可开关的实现}并在同一算例池上实测
（完整记录见支撑材料 \texttt{docs/optional\_changes.pdf}：11 项尝试中 3 项有效、
6 项为负优化、2 项无效）。结论如下。

\textbf{(1) 完成率损失——已解决。} 该损失此前被归因于"贴边且朝外的定向源在域内不可认证"，
但诊断显示真实原因在更上游：\\
(i) 末段三次观测点\textbf{近乎共线}（归航沿直线走），共线的方位测量无法约束距离，
最小包围圆半径因此停在 77 m，而末段清除图案的半径取 $0.55\sigma\approx42$ m，
图案点没有一个落进 20 m 清除半径；\\
(ii) 图案扫完后机器人停在\textbf{最后一个外圈点}上，该点可能位于定向源的\textbf{背光侧}，
随后的原地复测必然无声，循环空转直至预算耗尽。\\
两处修复分别是：\textbf{垂向截断测量}（图案失败时从垂直于视线的方向补一次观测，
$\sigma$ 由 77 m 降到 \textbf{1 m}）与\textbf{自适应沿射线站位}（第二站位偏移角随失败次数收缩，
最终回到射线上——由射线逼近引理保证必定受光）。在\textbf{从未参与调参的独立算例池}
（40 例）上，完成率由 $0.9945$ 升到 $\mathbf{1.0000}$（最差由 $0.8462$ 升到 $1.0000$），
平均时间同时下降 \textbf{38.2\%}。\textbf{真正的解不是域外采样，而是把定位做准、把站位放对。}

\textbf{(2) 路径规划的耦合——部分有效。} 把"清除顺序影响后续探测"具体化为两个可实现的耦合：
顺路清除（巡访站位时清除近旁的已知源）与自适应补测间隔（未解算频道多时缩短补测间隔）。
24 例对照：前者 $881.3\to882.2$ s（\textbf{无效}，说明任务排序中的优先权偏置已隐含实现了该耦合），
后者 $881.3\to\mathbf{852.1}$ s（$-3.3\%$，已默认启用）。因此本文的路径规划仍属启发式，
但"耦合未被利用"这一自评偏保守。

\textbf{(3) 误差独立同界——经检验无需建模。} 把示向度误差换成\textbf{空间相关}的平滑随机场
（波长 $1200\sim3000$ m，相关比例取 $0/50/80/100\%$）后，完成率与时间的差异落在运行波动范围内
（$a=100\%$ 时反而最好）。原因是结构性的：本文把每条读数当作\textbf{硬约束}
（真实方位落在 $\hat\theta\pm1^{\circ}$ 内），定位区域取全部约束之交、不确定度取最小包围圆半径，
该结论对\textbf{任意}满足逐条界的误差实现都成立，\textbf{与是否相关无关}。
相关性只影响"重测取平均"的收益，而本文重测的价值来自\textbf{交会角改善}而非误差平均。
（若改用卡尔曼滤波一类\textbf{统计}估计器，则必须先建模相关性。）

\textbf{(4) 真实时间耦合——本地环境判定无需建模，但在官方环境需复测。}
实测一次正式测试（问题四，799 次 HTTP 请求）的真实耗时为 5.18 s，
平均每请求 \textbf{6.48 ms}，相对 1200 s 上限有 \textbf{232 倍}余量；
即使计入仿真环境内部开销 19.63 s，余量仍有 61 倍。
但本结论只对本地模拟器成立：官方模拟器是 Windows 图形程序，其响应受界面刷新影响，
可能慢得多。本队无法获取官方模拟器（见第 10.1 节），故\textbf{无法实测}，
只能给出定量说法——官方每请求需比本地慢 230 倍以上才会触限。

\subsection{推广}

本文框架可直接推广到：(i) \textbf{多机器人协同}——按 Voronoi 划分站位并行搜索，终止判据变为“全体无信号点之并集覆盖目标域”；(ii) \textbf{移动干扰源}——估计改为带运动模型的滤波，覆盖判据改为时空覆盖；(iii) \textbf{TDOA/AOA 混合定位}——角度楔换成双曲线带，问题一算法推广为二次约束可行域枚举；(iv) \textbf{应急搜救}——替换“接收半径/频道”为“探测半径/目标类型”后三件套同样适用。

\begin{thebibliography}{99}
\bibitem{reg672} 中华人民共和国国务院, 中华人民共和国中央军事委员会. 中华人民共和国无线电管理条例（国令第 672 号）[Z]. 2016-11-11 修订, 2016-12-01 施行. \url{https://www.gov.cn/zhengce/content/2016-11/25/content_5137687.htm}
\bibitem{itu854} Recommendation ITU-R SM.854-3. Direction finding and location determination at monitoring stations[S]. ITU, 2011-09. \url{https://www.itu.int/rec/R-REC-SM.854/en}
\bibitem{itu2060} Recommendation ITU-R SM.2060-0. Test procedure for measuring direction finder accuracy[S]. ITU, 2014-11. \url{https://www.itu.int/rec/R-REC-SM.2060/en}
\bibitem{ituhdb} ITU-R. Handbook on Spectrum Monitoring (R-HDB-23-2011)[M]. Geneva: ITU, 2011. \url{https://www.itu.int/pub/R-HDB-23-2011}
\bibitem{gb34089} GB/T 34089-2017. VHF/UHF 无线电监测测向系统开场测试参数和测试方法[S]. 北京: 中国标准出版社, 2017.
\bibitem{yd3730} YD/T 3730-2020. VHF/UHF 无线电测向系统多径传播抗扰度测试程序[S]. 北京: 人民邮电出版社, 2020.
\bibitem{gb25003} GB/T 25003-2010. VHF/UHF 频段无线电监测站电磁环境保护要求和测试方法[S]. 北京: 中国标准出版社, 2010.
\bibitem{zhang2011} 张洪顺, 王磊. 无线电监测与测向定位[M]. 西安: 西安电子科技大学出版社, 2011. ISBN 978-7-5606-2654-3.
\bibitem{zhou2009} 周靖博, 黄艳. 做好无线电干扰源测向定位的几点经验[J]. 中国无线电, 2009(2).
\bibitem{huang2023} 黄裕文, 李怡恒, 程擎, 等. 机场终端区 GNSS 信号干扰源定位研究[J]. 现代雷达, 2023(12): 87-93.
\bibitem{torrieri1984} Torrieri D J. Statistical theory of passive location systems[J]. IEEE Transactions on Aerospace and Electronic Systems, 1984, AES-20(2): 183-198.
\bibitem{stansfield1947} Stansfield R G. Statistical theory of d.f. fixing[J]. Journal of the Institution of Electrical Engineers - Part IIIA: Radiocommunication, 1947, 94(15): 762-770.
\bibitem{xiu2005} 修建娟, 何友, 王国宏, 等. 测向交叉定位系统中的交会角研究[J]. 宇航学报, 2005, 26(3): 282-285.
\bibitem{sheng2016} 盛丹, 王国宏, 孙殿星. 存在系统误差下交叉定位系统最优交会角研究[J]. 系统工程与电子技术, 2016, 38(7): 1516-1523.
\bibitem{kershner1939} Kershner R. The number of circles covering a set[J]. American Journal of Mathematics, 1939, 61(3): 665-671.
\bibitem{gaspar2021} Gáspár Z, Tarnai T, Hincz K. Partial covering of a circle by 6 and 7 congruent circles[J]. Symmetry, 2021, 13(11): 2133.
\bibitem{koopman1946} Koopman B O. Search and Screening[M]. New York: Pergamon Press, 1980 (first edition: Operations Evaluation Group, 1946). ISBN 0080231357.
\bibitem{stone1975} Stone L D. Theory of Optimal Search[M]. New York: Academic Press, 1975. ISBN 0126724504.
\bibitem{stone2016} Stone L D, Royset J O, Washburn A R. Optimal Search for Moving Targets[M]. Springer, 2016. DOI 10.1007/978-3-319-26899-6.
\bibitem{chen2021} 陈建勇. 实用搜索理论[M]. 北京: 国防工业出版社, 2021. ISBN 978-7-118-12327-2.
\bibitem{choset2001} Choset H. Coverage for robotics - a survey of recent results[J]. Annals of Mathematics and Artificial Intelligence, 2001, 31(1-4): 113-126.
\bibitem{galceran2013} Galceran E, Carreras M. A survey on coverage path planning for robotics[J]. Robotics and Autonomous Systems, 2013, 61(12): 1258-1276.
\bibitem{ausiello2001} Ausiello G, Feuerstein E, Leonardi S, et al. Algorithms for the on-line travelling salesman[J]. Algorithmica, 2001, 29(4): 560-581.
\bibitem{afrati1986} Afrati F, Cosmadakis S, Papadimitriou C H, et al. The complexity of the travelling repairman problem[J]. RAIRO - Theoretical Informatics and Applications, 1986, 20(1): 79-87.
\bibitem{blum1994} Blum A, Chalasani P, Coppersmith D, et al. The minimum latency problem[C]//Proceedings of the 26th Annual ACM Symposium on Theory of Computing (STOC'94), 1994: 163-171.
\bibitem{psaraftis1995} Psaraftis H N. Dynamic vehicle routing: status and prospects[J]. Annals of Operations Research, 1995, 61(1): 143-164.
\bibitem{bersimas1991} Bertsimas D J, van Ryzin G. A stochastic and dynamic vehicle routing problem in the Euclidean plane[J]. Operations Research, 1991, 39(4): 601-615.
\bibitem{nardone1980} Nardone S C, Aidala V J. Necessary and sufficient observability conditions for bearings-only target motion analysis[R]. NUSC Technical Memorandum TM 80-2059 / DTIC AD-A102073, 1980.
\bibitem{nardone1981} Nardone S C, Aidala V J. Observability criteria for bearings-only target motion analysis[J]. IEEE Transactions on Aerospace and Electronic Systems, 1981, AES-17(2): 162-166.
\bibitem{lecadre1997} Le Cadre J E, Jauffret C. Discrete-time observability and estimability analysis for bearings-only target motion analysis[J]. IEEE Transactions on Aerospace and Electronic Systems, 1997, 33(1): 178-201.
\bibitem{payne1989} Payne A N. Observability problem for bearings-only tracking[J]. International Journal of Control, 1989, 49(3): 761-768.
\bibitem{cumcmthesis} LaTeX 工作室. CUMCMThesis: 全国大学生数学建模竞赛 LaTeX 论文模板[CP/OL]. \url{https://github.com/latexstudio/CUMCMThesis}
\end{thebibliography}
